% section: 差分から始まる離散的な微積分

\section{差分から始まる離散的な微積分}

\subsection{数列の変化を測る}

数列
\[
  a_0,a_1,a_2,\ldots
\]
の値だけを眺めていると,変化の規則が見えにくいことがある.
そこで,隣り合う二項の差を考える.
\[
  \Delta a_n=a_{n+1}-a_n
\]
これを数列の一階差分という.

例えば,
\[
  a_n=n^2
\]
ならば,
\[
  \Delta a_n=(n+1)^2-n^2=2n+1
\]
である.
さらに差分を取ると,
\[
  \Delta^2a_n
  =\Delta(2n+1)
  =(2(n+1)+1)-(2n+1)=2
\]
となる.

二次式が二回の差分で定数になった.
これは,二次関数を二回微分すると定数になることに対応している.
差分は,離散的な世界で変化の速さを測る操作なのである.

\subsection{差分を取ると次数が下がる理由}

多項式
\[
  p(n)=c_mn^m+c_{m-1}n^{m-1}+\cdots+c_0
\]
を考える.
二項定理から,
\[
  (n+1)^m-n^m
  =mn^{m-1}+\text{それより低い次数の項}
\]
である.
したがって,$\Delta p(n)$ の次数は,$p(n)$ の次数より一つ低い.

特に,$m$ 次多項式に $m$ 回差分を取ると定数になり,$m+1$ 回差分を取ると $0$ になる.
この事実から,数列が多項式で表されるかどうかを,差分表から推測できる.

\[
  \begin{array}{c|cccccc}
    n           & 0 & 1 & 2 & 3 & 4  & 5  \\ \hline
    a_n         & 0 & 1 & 4 & 9 & 16 & 25 \\
    \Delta a_n  & 1 & 3 & 5 & 7 & 9  &    \\
    \Delta^2a_n & 2 & 2 & 2 & 2 &    &
  \end{array}
\]

この表は,値の列よりも,変化の列の方が単純になることを示している.
漸化式を変形するという本編の発想も,実は「単純な変化を探す」という作業だった.

\subsection{差分演算子とシフト演算子}

差分の仕組みを,記号で整理してみよう.
数列を一つの対象と考え,シフト演算子 $E$ を
\[
  (Ea)_n=a_{n+1}
\]
で定める.
すると,
\[
  \Delta a_n=a_{n+1}-a_n=(E-1)a_n
\]
であるから,
\[
  \Delta=E-1
\]
と書ける.

この記号の利点は,差分を何回も取る操作を展開できることである.
\[
  \Delta^2=(E-1)^2=E^2-2E+1
\]
だから,
\[
  \Delta^2a_n=a_{n+2}-2a_{n+1}+a_n
\]
となる.
さらに,
\[
  \Delta^k=(E-1)^k
  =\sum_{j=0}^{k}(-1)^{k-j}\binom{k}{j}E^j
\]
より,
\[
  \Delta^ka_n
  =\sum_{j=0}^{k}(-1)^{k-j}\binom{k}{j}a_{n+j}
\]
を得る.

ここで現れた二項係数は偶然ではない.
差分は「一つ先へ進める操作」と「何もしない操作」の差であり,何度も差を取ると二項定理が現れるのである.

\subsection{差分の逆操作としての和}

微分の逆操作が積分であるように,差分の逆操作は和である.
例えば,
\[
  \Delta a_n=b_n
\]
ならば,
\[
  a_{n+1}-a_n=b_n
\]
なので,$0$ から $n-1$ まで足し合わせると,
\[
  a_n-a_0=\sum_{k=0}^{n-1}b_k
\]
となる.
従って,
\[
  a_n=a_0+\sum_{k=0}^{n-1}b_k
\]
である.

これは階差型漸化式の解法そのものである.
本編では解法として使ったが,ここでは「差分を元に戻す操作」として理解できる.

例えば,
\[
  a_{n+1}-a_n=2n+1,\qquad a_0=0
\]
ならば,
\[
  a_n=\sum_{k=0}^{n-1}(2k+1)
\]
であり,これは
\[
  a_n=n^2
\]
を与える.
平方数は,奇数を順番に足してできるという見方も得られた.

\subsection{和の公式を差分で設計する}

逆に,和
\[
  1^2+2^2+\cdots+n^2
\]
を求めたいとする.
公式を覚えていなくても,差分を利用して公式の形を設計できる.

和を $S_n$ とおくと,
\[
  S_n-S_{n-1}=n^2
\]
である.
右辺が二次式だから,$S_n$ は三次式になると予想する.
\[
  S_n=An^3+Bn^2+Cn+D
\]
と置けば,
\[
  \begin{aligned}
    S_n-S_{n-1}
     & =
    A\{n^3-(n-1)^3\}
    +B\{n^2-(n-1)^2\}
    +C   \\
     & =
    3An^2+(-3A+2B)n+(A-B+C)
  \end{aligned}
\]
である.
これが $n^2$ に一致するためには,
\[
  3A=1,\qquad -3A+2B=0,\qquad A-B+C=0
\]
である.
さらに $S_0=0$ より $D=0$ だから,
\[
  A=\frac13,\qquad B=\frac12,\qquad C=\frac16
\]
となる.
従って,
\[
  1^2+2^2+\cdots+n^2
  =\frac{n(n+1)(2n+1)}6
\]
が得られる.

ここで重要なのは,和の公式を暗記したことではない.
「差分を取ると元の被加数になるような式を探す」という,離散的な積分の発想を使ったことである.
